\chapter{Normativa y accesibilidad}

\section{Protección de datos}

El proyecto se orienta a una aplicación potencial en rehabilitación cognitiva, por lo que debe contemplar desde el diseño la protección de datos personales y, especialmente, de datos de salud. Durante el desarrollo descrito en esta memoria no se ha trabajado con datos reales de pacientes. Los registros analizados corresponden a vehículos autónomos simulados, métricas de entrenamiento y resultados técnicos.

En una futura validación clínica, cualquier uso con personas usuarias deberá ajustarse al Reglamento General de Protección de Datos y a la Ley Orgánica 3/2018 de Protección de Datos Personales y garantía de los derechos digitales~\cite{rgpd,lopdgdd}. Será necesario definir responsable del tratamiento, finalidad, base jurídica, minimización de datos, conservación, consentimiento informado cuando proceda, medidas de seudonimización y control de acceso. También deberá diferenciarse entre datos necesarios para la evaluación terapéutica y métricas técnicas internas del simulador.

\section{Seguridad y uso responsable}

El simulador no sustituye por sí mismo una evaluación clínica ni una autorización administrativa para conducir. Su función debe entenderse como apoyo a profesionales, entrenamiento controlado y herramienta de observación. Cualquier conclusión sobre capacidad real de conducción debería depender de especialistas y procedimientos establecidos.

Desde el punto de vista técnico, el sistema autónomo genera tráfico artificial. Por tanto, sus fallos no suponen riesgo físico directo durante el entrenamiento virtual, pero sí pueden afectar a la validez de la experiencia. Es importante que las sesiones con pacientes utilicen modelos suficientemente probados para evitar comportamientos incoherentes que confundan al usuario o introduzcan estímulos no deseados.

\section{Accesibilidad}

La accesibilidad del simulador debe analizarse en relación con las personas a las que se dirige, que pueden presentar alteraciones de atención, memoria, percepción o movilidad. En este TFG no se ha rediseñado la interfaz utilizada por el paciente, por lo que no puede afirmarse que el producto completo haya sido validado como accesible. La aportación realizada se concentra en el tráfico autónomo y afecta de forma indirecta a la carga cognitiva de cada sesión.

El número de vehículos, los puntos del mapa y la frecuencia de situaciones normativas deben poder ajustarse para que la dificultad aumente de forma progresiva. También es importante repetir un mismo escenario bajo condiciones comparables: si el tráfico cambia sin control, resulta difícil saber si una variación en el resultado procede del paciente o de una sesión más exigente. La separación entre entrenamiento y test de la IA y el registro de eventos sientan una base técnica para construir esas repeticiones, aunque todavía falta trasladar estos parámetros a una configuración manejable por el profesional.

La interacción debe contemplar periféricos adaptados, instrucciones breves y legibles y una presentación que no añada estímulos irrelevantes. El rendimiento gráfico y la duración de la prueba también deben vigilarse para reducir fatiga o mareo. Por último, los datos exportados no deberían mostrarse al especialista como un volcado técnico de CSV, sino como indicadores comprensibles y vinculados a maniobras concretas. Estos requisitos deberán revisarse con profesionales y usuarios de ASPAYM antes de utilizar el simulador en una validación clínica.
